\mode*
\subsection{More containers}
\label{containers}

After lists and dictionaries, an introductory course typically adds the
containers that answer more specialised questions: the tuple for a fixed
group of parts, the set for which distinct values occur, and the stack and the
queue for which element to take next.
Python writes them alike.
A tuple is indexed with the same brackets as a list, a dictionary is looked up
with the same brackets again, and \mintinline{python}{append} adds to a list
and to a queue alike.
\textcite{Johnson2020PythonMisconceptions} argue that this uniformity is itself
a source of misconceptions: in their analysis of Python novices,
\textcquote{Johnson2020PythonMisconceptions}{[o]verloading hides the complexity
of algorithms and data structures, often leading students to write code that
involves mutability, sharing, copying, side effects, coroutines, concurrency,
and lazy evaluation -- and none of those topics are accessible to students who
haven't yet mastered basic assignments, conditionals, and looping}.
When the surface is the same, what a container \emph{does} --- whether it can
change, what may serve as its key, whether it keeps duplicates, which end it
gives from --- is not visible in the program text, and must be discerned.

The literature on these containers at the introductory level is thin.
Documented difficulties with stacks and queues come from data-structures
courses (\enquote{CS2}), and they concern \emph{choosing} and
\emph{implementing} a structure more than predicting what one does
\autocite{Zingaro2018DataStructures}.
For the tuple's immutability, the requirements on a dictionary key and the
set's lack of duplicates and order, we found no study of novices'
misconceptions; the misconceptions analysed below are those our course
material is designed against, not yet documented in the literature.

\subsubsection{Tuples}

A tuple behaves like a list in every way a novice has used a list so far:
indexing, \mintinline{python}{len} and \mintinline{python}{for} work the same.
It is natural, then, to take the tuple to be a list with other brackets ---
or, having heard that it cannot be changed, a \enquote{frozen list}, a list
that happens to be locked.

The object of learning is \emph{what a tuple is}. The critical aspect is that
a tuple is immutable: once made, no element can be replaced, removed or added,
and an attempt to do so is refused with an error rather than silently ignored.
Since everything else about the two containers is alike, we hold the program
invariant and vary only the brackets.

\begin{description}
  \item[Contrast] The same pair in a list and in a tuple, with the same
    assignment to each. The list on the left changes; the tuple on the right
    raises a \mintinline{text}{TypeError} (the tuple \enquote{does not
    support item assignment}), so the difference lies in the container and
    not in the assignment.

    \hfill
    \begin{minipage}[t]{0.45\columnwidth}
      \begin{minted}[highlightlines={1}]{python}
        position = [58.4, 15.6]
        position[0] = 59.3
        print(position)
      \end{minted}
    \end{minipage}
    \hfill
    \begin{minipage}[t]{0.45\columnwidth}
      \begin{minted}[highlightlines={1}]{python}
        position = (58.4, 15.6)
        position[0] = 59.3
        print(position)
      \end{minted}
    \end{minipage}
    \newline
    The contrast alone teaches \emph{that} a tuple cannot be changed, which is
    compatible with the frozen-list view. What separates the tuple from a
    locked list is \emph{why} one would choose it: the number of parts is given
    by what the thing is --- a position has two, a colour three, a date three
    --- and the parts are not changed one at a time.

  \item[Generalisation] Keeping that criterion invariant, we vary the thing
    stored: a position, a colour and a date are tuples, whereas a guest list
    and a course's students grow and shrink, and are lists.
\end{description}

\subsubsection{Dictionary keys}

The value in a dictionary may be anything, including a list. The key may not:
it must be immutable, because the dictionary finds the entry from the key's
value, and a key whose value changed afterwards could not be found again. Two
beliefs predict otherwise. If a key is as unrestricted as a value, a list
serves as a key as well as a tuple does. If a dictionary is a list indexed by
strings, neither does.

The object of learning is \emph{what may serve as a key}, and it fuses two
aspects the student has met separately: immutability, from the tuple, and
lookup by value, from the dictionary.

\begin{description}
  \item[Fusion] We store a distance under a pair of cities, varying only
    whether the pair is a tuple or a list. The tuple key on the left works; the
    list key on the right raises \mintinline{text}{TypeError}, because a list
    is not hashable. The error is a \mintinline{text}{TypeError} and not a
    \mintinline{text}{KeyError}: the list is not a key that is missing, it
    cannot be a key at all.

    \hfill
    \begin{minipage}[t]{0.45\columnwidth}
      \begin{minted}[highlightlines={2}]{python}
        distance = {}
        distance[("Lund", "Ystad")] = 63
        print(distance)
      \end{minted}
    \end{minipage}
    \hfill
    \begin{minipage}[t]{0.45\columnwidth}
      \begin{minted}[highlightlines={2}]{python}
        distance = {}
        distance[["Lund", "Ystad"]] = 63
        print(distance)
      \end{minted}
    \end{minipage}
\end{description}

\subsubsection{Sets}

A set is written like a dictionary without values, and a set made from a list
looks like the list with its duplicates removed. That invites the view that a
set \emph{is} such a list --- one that happens to drop duplicates but otherwise
keeps its elements in order.

The object of learning is \emph{what a set is}. Two features are critical, and
both must be discerned: a set holds each value once, and it has no order. The
second is the one the list-without-duplicates view misses, and it shows in
comparison rather than in printing.

\begin{description}
  \item[Contrast] The same elements, in the same written order, in a list and
    in a set, varying only the container. The list has six elements, the set
    three; and the two comparisons, identical but for the brackets, give
    opposite answers.

    \hfill
    \begin{minipage}[t]{0.45\columnwidth}
      \begin{minted}{python}
        xs = [1, 2, 2, 3, 3, 3]
        print(len(xs))  # 6
        print([1, 2, 3] == [3, 2, 1])
        # False
      \end{minted}
    \end{minipage}
    \hfill
    \begin{minipage}[t]{0.45\columnwidth}
      \begin{minted}{python}
        xs = {1, 2, 2, 3, 3, 3}
        print(len(xs))  # 3
        print({1, 2, 3} == {3, 2, 1})
        # True
      \end{minted}
    \end{minipage}
    \newline
    Two sets are equal when they hold the same values; two lists are equal
    only when they hold the same values in the same order. That the set has
    no positions also means that \mintinline{python}{xs[0]} has nothing to
    refer to, and that a set of strings may print its elements in a different
    order from one run to the next.
\end{description}

\subsubsection{Stacks and queues}

A stack and a queue answer the same question --- which element is taken next?
--- with opposite answers: the stack gives back the element added last, the
queue the element added first. In Python the two are written almost alike, a
list with \mintinline{python}{append} and \mintinline{python}{pop} for the
stack and a \mintinline{python}{collections.deque} with
\mintinline{python}{append} and \mintinline{python}{popleft} for the queue, so
the order must be discerned from the method, not from the shape of the code.

The order itself is not where students in data-structures courses struggle
most. \textcite[p.~170]{Zingaro2018DataStructures} report an earlier finding
that students did best on questions about stack and queue interfaces, and in
their own study only seven per cent chose a structure without the
last-in-first-out property for an undo function
\autocite[p.~173]{Zingaro2018DataStructures}. At the introductory level the
risk is rather that one order is carried over to both containers.

The object of learning is \emph{which element a stack and a queue give next};
the critical aspect is the end the element is taken from.

\begin{description}
  \item[Contrast] The same elements added in the same order, varying only the
    container and therefore the end. The names are deliberately neutral:
    \mintinline{python}{stack} and \mintinline{python}{queue} would say which
    order to expect.

    \hfill
    \begin{minipage}[t]{0.45\columnwidth}
      \begin{minted}[highlightlines={5}]{python}
        things = []
        for x in [1, 2, 3]:
            things.append(x)
        while things:
            print(things.pop())
        # 3 2 1
      \end{minted}
    \end{minipage}
    \hfill
    \begin{minipage}[t]{0.45\columnwidth}
      \begin{minted}[highlightlines={6}]{python}
        import collections
        things = collections.deque()
        for x in [1, 2, 3]:
            things.append(x)
        while things:
            print(things.popleft())
        # 1 2 3
      \end{minted}
    \end{minipage}
\end{description}

\subsubsection{Choosing a container}

The analyses above come together when a container is to be chosen. The choice
is governed by the operations the program needs --- look a value up by a name,
ask whether something occurs, keep an order, keep a fixed group of parts, take
the newest or the oldest next --- and not by what the data is about: the same
students may be kept in a list, a set or as dictionary keys, depending on
whether the program enumerates them in order, asks whether someone is
enrolled, or looks something up for each of them.
Among the students \textcite[p.~173]{Zingaro2018DataStructures} asked to
choose an implementation for a generalised array, the twelve per cent who
included a slow structure justified the choice by what \emph{could} be built
with it rather than by the cost of the operations needed; in one student's
words,
\textcquote{Zingaro2018DataStructures}{[a]ll three data structures should work
as long as they arrange the elements by subscript}.

\begin{description}
  \item[Generalisation] Keeping the criterion invariant --- the operations
    the program needs --- we vary the data: the registered students (asked
    whether someone is enrolled: a set), their grades (looked up per student:
    a dictionary), the order the solutions arrived in (graded first come,
    first served: a queue), the date of a submission (a fixed group of parts:
    a tuple).
\end{description}

Because several containers can be defended for the same data, this aspect
does not lend itself to a single keyed answer; the diagnostic quizzes
(\cref{app:diagnostics-containers}) therefore probe the four aspects it is
built from rather than the choice itself.
